\section{ELBO 在扩散模型中的推导}

\subsection{扩散模型的 ELBO}

与 VAE 类似，扩散模型的训练目标是最大化数据的对数似然的下界。对于层级化的 VAE，ELBO 为：

$$\log p_\theta(x_0) \geq \mathbb{E}_{q(x_{1:T}|x_0)} \left[ \log \frac{p_\theta(x_{0:T})}{q(x_{1:T}|x_0)} \right]$$

\subsection{ELBO 的第一种分解}

将 ELBO 展开，可以得到三类项：

$$-\text{ELBO} = \underbrace{-\mathbb{E}_{q(x_1|x_0)}[\log p_\theta(x_0|x_1)]}_{\text{重构项 } L_0} + \underbrace{D_{\text{KL}}(q(x_T|x_0) \| p(x_T))}_{\text{先验匹配项 } L_T} + \underbrace{\sum_{t=2}^{T} L_{t-1}}_{\text{一致性项}}$$

其中一致性项为：

$$L_{t-1} = D_{\text{KL}}(q(x_{t-1}|x_t, x_0) \| p_\theta(x_{t-1}|x_t))$$

\begin{importantbox}{ELBO 三项的含义}
\begin{enumerate}
    \item \textbf{重构项 $L_0$}：衡量从 $x_1$ 重构 $x_0$ 的质量，与 VAE 的重构损失类似。在实践中，$x_1$ 与 $x_0$ 非常接近（仅加了很少的噪声），因此这一项通常很小。
    \item \textbf{先验匹配项 $L_T$}：衡量前向过程的终点 $q(x_T|x_0)$ 与先验 $p(x_T) = \mathcal{N}(0, I)$ 的距离。由于前向过程是固定的，且当 $T$ 足够大时 $q(x_T|x_0) \approx \mathcal{N}(0,I)$，这一项接近零且不依赖于 $\theta$。
    \item \textbf{一致性项 $L_{t-1}$}：这是\textbf{核心训练项}，衡量学习的逆过程 $p_\theta(x_{t-1}|x_t)$ 与真实的``一步逆过程'' $q(x_{t-1}|x_t, x_0)$ 之间的差距。
\end{enumerate}
\end{importantbox}

\subsection{马尔可夫假设的简化}

在马尔可夫假设下，ELBO 的分解可以进一步简化。关键是以下等式成立：

$$q(x_t | x_{t-1}, x_0) = q(x_t | x_{t-1})$$

这是因为在马尔可夫链中，给定 $x_{t-1}$ 后，$x_t$ 与 $x_0$ 条件独立。这一性质使得一致性项中出现的后验条件分布 $q(x_{t-1}|x_t, x_0)$ 可以通过贝叶斯公式计算（这将在后续讲座中详细推导）。

\begin{knowledgebox}{条件后验分布的可计算性}
虽然 $q(x_{t-1}|x_t)$ 本身是难以计算的（需要对所有可能的 $x_0$ 做积分），但条件版本 $q(x_{t-1}|x_t, x_0)$（同时给定 $x_t$ 和 $x_0$）却有闭式高斯解。这是因为：
\begin{itemize}
    \item $q(x_t|x_{t-1})$ 是高斯（前向转移）
    \item $q(x_{t-1}|x_0)$ 是高斯（前向过程闭式解）
    \item 高斯的条件仍然是高斯
\end{itemize}
这使得一致性项可以精确计算为两个高斯之间的 KL 散度。
\end{knowledgebox}

\subsection{先验匹配项的自然满足}

前向过程的闭式解告诉我们：

$$q(x_T | x_0) = \mathcal{N}(\sqrt{\bar{\alpha}_T} x_0, (1-\bar{\alpha}_T) I)$$

当 $T$ 足够大时，$\bar{\alpha}_T = \prod_{t=1}^T \alpha_t \to 0$（因为每个 $\alpha_t < 1$），所以：

$$q(x_T | x_0) \to \mathcal{N}(0, I) = p(x_T)$$

因此 $L_T = D_{\text{KL}}(q(x_T|x_0) \| p(x_T)) \to 0$。

\begin{importantbox}{前向过程的设计目标}
前向过程被精心设计为：通过逐步添加噪声并缩放信号，使得经过足够多步后，\textbf{无论初始数据 $x_0$ 是什么}，最终的 $x_T$ 都接近标准高斯分布。这是非递减 $\beta_t$ 和 $T$ 足够大的共同要求。前向过程的设计使得 ELBO 中的先验匹配项自动被满足，无需额外优化。
\end{importantbox}

\subsection{本章小结}

扩散模型的 ELBO 可以分解为重构项、先验匹配项和一致性项。前向过程的巧妙设计使先验匹配项自动趋近零，重构项通常也很小。因此，训练的核心在于最小化一致性项——即让逆过程 $p_\theta$ 逐步匹配真实的条件后验 $q(x_{t-1}|x_t, x_0)$。


